% Generated draft: requires mathematical and source review.
\documentclass[11pt]{article}
\usepackage[margin=25mm]{geometry}
\usepackage{fontspec}
\setmainfont{Latin Modern Roman}
\usepackage{amsmath,amssymb,mathtools,unicode-math}
\setmathfont{Latin Modern Math}
\usepackage{enumitem,needspace,xurl,hyperref}
\setlength{\parindent}{0pt}
\setlength{\parskip}{5pt}
\setlength{\emergencystretch}{3em}
\newcommand{\N}{\mathbb{N}}
\newcommand{\Z}{\mathbb{Z}}
\newcommand{\Q}{\mathbb{Q}}
\newcommand{\R}{\mathbb{R}}
\newcommand{\C}{\mathbb{C}}
\newcommand{\F}{\mathbb{F}}
\newcommand{\A}{\mathbb{A}}
\newcommand{\PP}{\mathbb{P}}
\newcommand{\E}{\mathbb{E}}
\newcommand{\eps}{\varepsilon}
\newcommand{\dd}{\,\mathrm{d}}
\begin{document}
\section*{UnsolvedMath ID: 30000492}
\textbf{Source problem code:} catalogue entry
\subsection{Definitions and mathematical statement}
\textbf{Settledness and Markov questions.} Is every irreducible quadratic polynomial \(f\in\F_p[x]\) settled? More generally, does the large-\(n\) distribution of critical-orbit labels of irreducible factors of its iterates follow the long-term behavior of the appropriate Markov process? Both questions are retained; the process and label rules are not specified in the record.\par Write \(f^{\circ n}\) for the \(n\)-fold compositional iterate, of degree \(2^n\), not the ordinary power \(f(x)^n\). In the usual relative stability convention, an irreducible polynomial \(g\) is \(f\)-stable if \(g\circ f^{\circ m}\) is irreducible over \(\F_p\) for every \(m\ge0\). If \(f^{\circ n}=c_n\prod_i g_{n,i}^{e_{n,i}}\) is its irreducible factorization, the source's degree-proportion definition of settledness becomes\[\lim_{n\to\infty}\frac{\sum_{i:\ g_{n,i}\text{ is }f\text{-stable}}e_{n,i}\deg g_{n,i}}{2^n}=1.\] This weights factors by degree, with multiplicities included if present, rather than counting each factor equally. The exact source convention for \(f\)-stable is not quoted separately, so the conventional future-composition definition above needs checking.\par For odd \(p\), a quadratic \(f(x)=ax^2+bx+c\), \(a\ne0\), has finite critical point \(-b/(2a)\) and a finite forward orbit over \(\F_p\). Common factor types use evaluations along that orbit, for example their square/nonsquare characters, but the source does not fix the orbit's starting index, repeats, labels, or handling of zero evaluations. These examples do not define its actual Markov chain. Characteristic two needs a different critical-point convention; an odd-prime restriction is not present and is not silently inserted into the literal universal claim.\par The Markov assertion requires a state space, transition probabilities, initial distribution, and a precise empirical distribution on factors. Degree-weighted and uniformly factor-counted distributions differ. It also requires a convergence or asymptotic comparison criterion: following a chain's long-term behavior does not automatically mean convergence to a unique stationary law. The class meant by \emph{more generally}, and whether the process is a theorem about factorization or a heuristic model, are absent. No transition rule or stationary measure is invented.\par The supplied triage checked 2026-08-21 labels the general picture open. It reports partial evidence from settled elements and special postcritically finite cases, while saying these do not establish the general factor-distribution assertion. Exact theorem hypotheses and their relation to both questions are not reproduced. This dated assessment requires independent verification.
\subsection{Short English statement}
The first question asks whether stable branches occupy asymptotically all the degree of compositional iterates of an irreducible quadratic. The second asks for a critical-orbit Markov description of factor types. Stability, labeling, characteristic-two treatment, and the process itself need source definitions.
\subsection{Sources}
\begin{itemize}\item Local UnsolvedMath registry, permanent ID 30000492, code catalogue entry; \emph{Settled Quadratic Polynomials and Markov Factorization Models}, Oberwolfach Reports Open Problems.\item \emph{Pro-\(p\) Extensions of Global Fields and pro-\(p\) Groups (2007)}, Oberwolfach Reports, DOI 10.4171/owr/2006/25, pp. 1463--1536. \url{https://doi.org/10.4171/owr/2006/25}. Proposed year and year tag: 2006; no proposer named. The source's escaped title delimiters are rendered as mathematics.\item Supplied literature label \emph{original OWR report}: \url{https://doi.org/10.4171/owr/2006/25}.\item Supplied literature label \emph{Boston--Jones arboreal representations}: \url{https://people.carleton.edu/~rfjones/Preprints/arborealrepnsfinal.pdf}.\item Supplied literature label \emph{recent special-case result}: \url{https://arxiv.org/abs/2604.04524}.\item Literature triage checked 2026-08-21, status \texttt{open}; record updated 2026-09-16. The \texttt{corrected\_verified} field reports recovering both questions and settledness from the official report, not independently checked here.\end{itemize}
\end{document}
